%!TEX root = ../../../exa-ma-d7.1.tex
\section{Software: Feel++}
\label{sec:WP7:Feel++:software}

\begin{table}[!ht]
    \centering
    { \setlength{\parindent}{0pt}
    \def\arraystretch{1.25}
    \arrayrulecolor{numpexgray}
    {\fontsize{9}{11}\selectfont
    \begin{tabular}{!{\color{numpexgray}\vrule}p{.4\textwidth}!{\color{numpexgray}\vrule}p{.6\textwidth}!{\color{numpexgray}\vrule}}
        \rowcolor{numpexgray}{\rule{0pt}{2.5ex}\color{white}\bf Field} & {\rule{0pt}{2.5ex}\color{white}\bf Details} \\
        \rowcolor{white}\textbf{Consortium} & \begin{tabular}{l}
Feel++ Consortium\\
\end{tabular} \\
        \rowcolor{numpexlightergray}\textbf{Exa-MA Partners} & \begin{tabular}{l}
CNRS\\
Inria Grenoble\\
Unistra\\
\end{tabular} \\
        \rowcolor{white}\textbf{Contact Emails} & \begin{tabular}{l}
christophe.prudhomme@cemosis.fr\\
vincent.chabannes@cemosis.fr\\
\end{tabular} \\
        \rowcolor{numpexlightergray}\textbf{Supported Architectures} & \begin{tabular}{l}
CPU Only\\
\end{tabular} \\
        \rowcolor{white}\textbf{Repository} & \href{https://github.com/feelpp/feelpp}{https://github.com/feelpp/feelpp} \\
        \rowcolor{numpexlightergray}\textbf{License} & \begin{tabular}{l}
OSS:: GPL v*\\
OSS:: LGPL v*\\
\end{tabular} \\
        \bottomrule
    \end{tabular}
    }}
    \caption{WP7: Feel++ Information}
\end{table}

\subsection{Software Overview}
\label{sec:WP7:Feel++:summary}

Feel++ is an open-source finite-element platform written in modern C++ and co-developed by CNRS, Inria, and the University of
Strasbourg. In WP7 it serves as a flagship codebase to exercise the application-framework taxonomy: the team maintains a
catalogue of mini-apps, reduced-basis demonstrators, and full digital twins that are packaged with shared GitHub Actions,
container images, and benchmarking recipes. The collaboration also channels upstream developments from WP2 (model-order
reduction) and WP6 (uncertainty quantification) into reproducible software artefacts.

\begin{table}[!ht]
    \centering
    {
        \setlength{\parindent}{0pt}
        \def\arraystretch{1.25}
        \arrayrulecolor{numpexgray}
        {
            \fontsize{9}{11}\selectfont
            \begin{tabular}{!{\color{numpexgray}\vrule}p{.25\linewidth}!{\color{numpexgray}\vrule}p{.6885\linewidth}!{\color{numpexgray}\vrule}}

    \rowcolor{numpexgray}{\rule{0pt}{2.5ex}\color{white}\bf Features} &  {\rule{0pt}{2.5ex}\color{white}\bf Short Description }\\

\rowcolor{white}    Health & Cardiovascular and respiratory digital twin workflows validated on HPC clusters (pulse wave propagation, bronchial flows). \\
\rowcolor{numpexlightergray}    Environment & Coastal and atmospheric transport studies leveraging high-order FEM and adaptive mesh refinement. \\
\rowcolor{white}    Energy & Thermal-hydraulics and magneto-hydrodynamics demonstrators linked to WP5 optimization campaigns. \\
\rowcolor{numpexlightergray}    Physics & Generic PDE engine supporting mixed formulations, reduced bases, and uncertainty quantification. \\
\rowcolor{white}    mini-apps & Parametric Feel++ mini-app catalogue (geometry scans, ROM kernels) integrated with WP7 ReFrame suites. \\
\end{tabular}
        }
    }
    \caption{WP7: Feel++ Features}
\end{table}


\subsection{Parallel Capabilities}
\label{sec:WP7:Feel++:performances}


Feel++ exposes distributed-memory parallelism via MPI and leverages PETSc/Hypre back-ends for sparse linear algebra, with
optional GPU acceleration through CUDA-enabled PETSc builds. Thread-level parallelism (OpenMP and TBB) is used inside assembly
kernels generated by the Feel++ DSL. The Exa-MA deployment targets the Jean Zay CPU partition, the IDRIS Joliot-Curie AMD nodes,
and the CINES Adastra hybrid system, relying on WP7-provided Spack recipes to ensure reproducible builds. Recent benchmarks
demonstrate near-linear scaling up to 1024 cores on the ``feelpp-io'' mini-app and sustained 70\% efficiency for the arterial
wall demonstrator when moving from 8 to 64 GPU-enabled ranks. Feel++ integrates seamlessly with WP7 automation by packaging its
ReFrame test suite and container artefacts alongside shared documentation templates.


\subsection{Initial Performance Metrics}
\label{sec:WP7:Feel++:metrics}

Baseline WP7 measurements target the ``app-feelpp-discr-1'' mini-app (elliptic problem with geometric parametrization) and the
arterial flow demonstrator. Input geometries and parameter sweeps are shared through the WP7 Feel++ Zenodo community (deposition scheduled Q1~2026), with automated
ReFrame recipes generating reference outputs and timing logs. On Jean Zay's Rome partition the mini-app reaches 0.14~s per solve
at 256 cores with 94\% weak-scaling efficiency; the digital twin scenario runs in 38~s per cardiac cycle on 512 cores while
sustaining I/O throughput above 12~GB/s thanks to the WP7 asynchronous HDF5 pipeline. Identified bottlenecks include PETSc
preconditioner setup costs for highly anisotropic meshes and initial container start-up time on Adastra's Cray OS, both of which
are being profiled in collaboration with WP3.

\subsubsection{Benchmark \#1}
\begin{itemize}
    \item \textbf{Description:} ``app-feelpp-discr-1'' weak-scaling run on Jean Zay (up to 512 Rome cores) solving a parametrised Poisson problem with $10^8$ unknowns.
    \item \textbf{Benchmarking Tools Used:} ReFrame orchestrates executions, Score-P captures MPI timelines, and the WP7 energy plug-in collects node-level power traces.
    \item \textbf{Input/Output Dataset Description:}
        \begin{itemize}
            \item \textbf{Input Data:} Parametrised geometry pack (1.2~GB HDF5) distributed via the WP7 Feel++ Zenodo community archive.
            \item \textbf{Output Data:} Solution snapshots and PETSc performance logs stored alongside the ReFrame metadata (NetCDF + JSON).
            \item \textbf{Data Repository:} Zenodo community ``Feel++ Exa-MA benchmarks'' maintained jointly by WP7 and Cemosis.
        \end{itemize}
    \item \textbf{Results Summary:} 94\% weak-scaling efficiency between 64 and 512 cores; average energy-to-solution 27\% lower than pre-container baseline due to improved process pinning.
    \item \textbf{Challenges Identified:} Residual load imbalance from mesh partitioning manifests as 8\% spread in time-to-solution; WP7 is testing ParMETIS and Zoltan repartitioning strategies.
\end{itemize}

\subsection{12-Month Roadmap}
\label{sec:WP7:Feel++:roadmap}

\begin{itemize}
    \item \textbf{Data improvements:} Publish curated mesh/parameter bundles for cardiovascular use-cases with automatic metadata capture from Feel++ scripts (Q1~2026).
    \item \textbf{Methodology application:} Extend the WP7 CI to run nightly reduced-basis regression tests on Adastra's hybrid partition, including energy traces (Q2~2026).
    \item \textbf{Results retention:} Deploy a Grafana dashboard fed by WP7 ReFrame exports to visualise trend lines for convergence, time-to-solution, and energy per run (Q3~2026).
\end{itemize}